\chapter*{선행 사례에서 배운 것: 문제의식과 기준의 출처}
\addcontentsline{toc}{chapter}{선행 사례에서 배운 것}
\markboth{선행 사례에서 배운 것}{}

이 장은 이전 기수의 우수 사례와 문헌에서 무엇을 보고, 어디가 문제라고 판단했으며, 무엇을 가져와 어떻게 고쳤는지를 출처와 함께 적는다. 이 프로젝트의 기준(수익률 정의, 학습 방식, 비교 방식)은 대부분 여기서 출발했다.

\section*{출처 목록}
\begin{table}[H]\centering\small
\begin{tabularx}{\linewidth}{@{}p{12mm}p{48mm}L@{}}\toprule
기호 & 출처 & 이 보고서에서 쓴 곳\\\midrule
{}[P1] & 이전 기수 우승 보고서(이하 ``우승 보고서'') & 수익률 정의, 승인 규칙, 결론 비교\\
{}[P2] & 지난 우수 보고서의 부도율 표 & 부도율 정의\\
{}[P3] & 지난 발표: 4조 & 시점 구분(누구의 결정인가)\\
{}[P4] & 지난 발표: 9기 5조 & 부도 정의, 누수, 손실 회피의 계산\\
{}[P5] & 지난 발표: 2025년 발표 & 결측 처리, 분할, Sharpe 임계값\\
{}[P6] & 수업 과제 정의(과제식 Sharpe: 부도 = 0, 무위험 5\%, 거절 = 초과수익 0) & 평가 지표\\
{}[L1] & Serrano-Cinca \& Guti\'errez-Nieto (2016), 수익 기반 점수(profit scoring) & 선별 곡선의 선행 연구\\
{}[L2] & Verbraken 외 (2014), EMP(기대 최대 이익) 척도 & 선별 곡선의 선행 연구\\
{}[L3] & Li, Li, Bellotti, Yao (2023), \emph{European Journal of Operational Research}, 부도와 조기상환의 경쟁 위험 모형 & 기대수익 점수의 한계와 개선 방향\\\bottomrule
\end{tabularx}
\caption{[P]는 이전 기수 자료, [L]은 문헌. 이전 기수 자료의 쪽 번호는 원문 대조 후 보완한다. 재현 코드와 결과는 \ft{notebooks/evidence/}(\ft{paper_check.py}, \ft{RESULTS.md}), 판단 기록은 \ft{docs/decisions.md} 1, 3, 4, 27--28단계.}
\end{table}

\section*{왜 이 문제의식을 가졌나}
처음 출발점은 [P1]의 결론이었다. 우승 보고서는 ``LC 대출을 전부 사면 국채보다 못하다''고 보고, 그래서 골라서 승인해야 한다는 논리를 폈다. 우리가 같은 데이터를 다시 재 보니, 이 결론은 데이터가 아니라 \textbf{재는 방법}에서 나왔다. 수익률을 어떻게 정의하느냐에 따라 같은 2013년 대출의 수익률이 2.8\%(단순 연환산)에서 8.9\%(현금흐름 IRR)까지 달라졌고, ``국채를 이기는가''의 답이 뒤집혔다(2장). 그래서 이 프로젝트의 첫 질문은 ``어떤 모델이 더 잘 맞히나''가 아니라 \textbf{``성과를 정직하게 재는 방법은 무엇인가''}가 되었다. 요약의 ``정직한 성적표''는 여기서 나온 말이다.

\section*{아이디어별: 무엇을 보고, 무엇이 문제였고, 어떻게 고쳤나}
\begin{table}[H]\centering\footnotesize
\begin{tabularx}{\linewidth}{@{}p{20mm}p{12mm}LLp{12mm}@{}}\toprule
주제 & 출처 & 본 것과 문제점 & 우리가 확인하고 바꾼 것 & 위치\\\midrule
수익률 정의 & [P1] & 받은 총액으로 단순 연환산. 매달 돌아오는 원리금을 무시해, 약정대로 다 갚아도 평균 6.6\%(약정 금리 13.2\%)로 나온다 & 월별 현금흐름을 복원한 코호트 IRR, 수수료 1\% 차감. 2013년 2.8\% \ensuremath{\to} 8.9\%. 전부 승인 5.73\% vs 국채 0.87\%, 48개월 모두 국채보다 높음 & 2장, 7장\\
평가 표본 & [P1] & 최근 연도의 끝난 대출만으로 평가. 일찍 끝나는 부도가 과하게 섞임 & 99.99\% 종결된 2013--16년 36개월만 주 표본. 종결 편향을 따로 실험(26단계) & 3장\\
결론 재현 & [P1] & ``전부 승인 < 국채'' & 같은 설정으로 그 결론이 나오는 경로를 재현(Sharpe $-0.113$ vs 보고서 $-0.154$, 완전 재현은 못 함, 잠정) & 부록 C\\
무위험 수익률 & [P6] & 고정 5\%. 이 기간 실제 국채 평균(약 1.4\%)보다 크게 높음 & 발행월과 만기를 맞춘 국채(GS3, GS5). 회사채(BAA) 등 허들 4단계 민감도 & 2장, 7장\\
Sharpe 정의 & [P6], [P5] & 거절을 초과수익 0으로 넣는 대출 건별 Sharpe는 승인비율이 줄면 기계적으로 변한다(항등식 $\sqrt{a}\,m/\sqrt{s^2+(1-a)m^2}$). [P5]에서 Sharpe 최고 임계값이 ``극소수 승인''으로 나온 것도 같은 효과 & 같은 승인비율의 무작위와 비교하는 선별 가치를 주 판정으로, 발행월 포트폴리오 Sharpe를 보조로 & 6장, 9장\\
데이터 분할 & [P1], [P5] & 무작위 분할. 미래 대출로 과거를 맞혀 성적이 부풀려짐 & 시점 학습. AUC 0.73 \ensuremath{\to} 0.62--0.68(실제로 가능한 점수) & 3장\\
결측 처리 & [P5] & 결측 행을 모두 삭제 & 결측 유형 분류(구조적, 사건 없음, 시기, MNAR). LightGBM은 결측 그대로, 로지스틱은 가용성 규칙 & 4장, 기술 부록\\
누수 & [P4] & 받은 총액(\ft{total_pymnt})을 입력에 넣음 & ``승인 직전에 알 수 있는가''로 누수 목록을 명시하고 테스트로 막음 & 3장\\
시점 구분 & [P3] & \textbf{영감:} 신청 \ensuremath{\to} LC 심사 \ensuremath{\to} 투자 모집 \ensuremath{\to} 상환 \ensuremath{\to} 사후로 ``누구의 결정인가''를 나눔 & 투자자는 LC가 금리, 등급을 매긴 뒤 고르므로 금리, 등급, 월 상환액은 사용 가능. 같은 기준으로 \ft{funded_amnt_inv}를 약한 누수 후보로 찾아 빼 봤으나 결과 같음(+0.01\%p) & 8장\\
부도 정의와 손실 계산 & [P4] & ``IRR < 국채''를 부도로 정의하고, 거절한 부도 대출의 피한 손실을 수익처럼 더함 & 부도는 상태(상각, 부도)로 정의. 거절한 돈은 국채에 두고(초과수익 0) 실제로 받은 돈만 채점 & 1장, 9장\\
부도율 표기 & [P2] & ``부도율 24\%, 33\%''는 부도 \ensuremath{\div} 정상 비율 & 실제 부도율 19.3\%, 25.1\%로 정정. 지표마다 분모를 명시 & 8장\\
승인 규칙 & [P1] & \textbf{영감:} 기대수익과 위험의 비율로 십분위를 나눠 승인. 다만 $\sigma$가 기대수익의 함수라 위험 정보를 더하지 못함 & 기대수익 점수 $s=(1-p)r-p\cdot\mathrm{LGD}-$국채(LGD는 그 시점까지 끝난 부도의 평균)를 시점 학습 $p$로 계산. 최종 제안의 기준이 됨 & 5장, 9장\\
선별 곡선 & [L1], [L2] & 처음에는 ``새 방법''으로 생각 & 선행 연구가 있음을 확인하고 정정. 우리 기여는 시점 학습, 코호트 IRR, 같은 비율 무작위 비교를 합친 \textbf{채점 설계} & 8장\\
조기상환 & [L3] & 부도와 조기상환을 함께 모형화 & 우리 $s$는 조기상환을 무시한 1년 근사라는 한계로 적고, 개선 방향으로 남김 & 10장\\\bottomrule
\end{tabularx}
\caption{선행 사례와 문헌에서 가져온 것과 고친 것. 숫자는 train 결과.}
\end{table}

\section*{정리: 무엇을 이어받고 무엇을 바꿨나}
\begin{itemize}[leftmargin=*]
\item \textbf{이어받은 것:} 기대수익으로 대출을 고른다는 발상([P1]), 시점별로 ``누구의 결정인가''를 나누는 틀([P3]), IRR과 만기 맞춘 국채를 쓴다는 방향([P5]).
\item \textbf{바꾼 것:} 수익률은 실제 현금흐름 IRR로, 무위험 수익률은 그 달 그 만기의 국채로, 학습은 그때 알던 정보로만, 비교는 같은 비율의 무작위와, 판정 기준은 test 전에 고정.
\item \textbf{그 결과 달라진 결론:} ``전부 승인은 국채보다 못하니 골라야 한다'' \ensuremath{\to} ``전부 사도 국채보다 나았고, 모델로 고르면 같은 양을 무작위로 사는 것보다 더 번다''(핵심 정리, 9장).
\end{itemize}
\begin{warnbox}[한계]
이전 기수 자료는 공개된 설명대로 다시 구현해 비교했다(원 보고서는 XGBoost와 10\% 표본, 우리는 LightGBM과 전체 표본). 그래서 수치를 완전히 재현하지는 못했다. 원문의 쪽 번호와 정확한 인용은 원문 대조 후 보완한다.
\end{warnbox}
